\labsection{2}{Audit missingness, categories, and validity}{sec:lab02}

\begin{labproject}{Data-quality audit: define rules before cleaning}
\textbf{Coverage.} Chapter~\ref{ch:quality}.\\
\textbf{Goal.} You should be able to distinguish missingness, invalid values, inconsistent categories, duplicate records, and structural absence; then choose a defensible treatment based on variable meaning and analytical purpose.\\
\textbf{Datasets.} \texttt{iris\_dataset.csv} and \texttt{titanic\_dataset.csv}; Ames Housing and Red Wine Quality are used as extensions.\\
\textbf{Scenario.} You have received data from multiple sources and must decide whether it is fit for analysis. The objective is not to make the data ``look clean'', but to identify defects, document rules, and preserve an auditable record of every transformation.

\textbf{Core audit.}
\begin{enumerate}
  \item \textbf{Create a quality profile.} For both primary datasets, report dimensions, variable classes, key uniqueness, and missing count/percentage for every variable.
  \item \textbf{Define validity rules.} Write rules before testing them: positive flower measurements; valid Titanic class and survival codes; non-negative fares; and plausible identifier uniqueness.
  \item \textbf{Standardise categories safely.} Inspect raw Iris species labels, create a new cleaned label, and verify that accidental spelling/format variants collapse to the intended categories. Do not overwrite the raw field.
  \item \textbf{Investigate missingness.} For Titanic, compare the concentration of missing values across \texttt{Age}, \texttt{Cabin}, and \texttt{Embarked}. Explain why one treatment is not appropriate for all three variables.
  \item \textbf{Demonstrate transparent imputation.} Create \texttt{Age\_imputed} using a documented group-median rule while preserving \texttt{Age}. Report how many values were changed and which grouping variables determined the replacement.
  \item \textbf{Check consequences.} Compare the age distribution before and after imputation. State one benefit and one distortion that the chosen rule can introduce.
\end{enumerate}

\begin{tryit}
\textbf{Extension A---structural missingness.} In \texttt{ames\_housing\_dataset.csv}, inspect variables with the highest missingness. Identify at least three fields where missing may mean ``feature absent'' rather than ``measurement unknown''. Explain how that semantic distinction changes preprocessing.

\textbf{Extension B---duplicate semantics.} In \texttt{red\_wine\_quality\_dataset.csv}, count exact duplicate rows. Before removing any, define the intended observational unit and explain why identical measurements do not automatically prove that a row is erroneous.
\end{tryit}

\tcbbreak
\lstinputlisting[style=utpr]{all_experiments/lab02_data_quality.R}
\begin{verification}
\textbf{Required evidence.} Submit a data-quality log with each issue written as \emph{rule $\rightarrow$ evidence $\rightarrow$ action $\rightarrow$ consequence}. Include at least one issue you deliberately leave unchanged because the evidence is insufficient to justify a transformation.

\textbf{Professional standard.} A cleaning step is acceptable only when another analyst can determine what changed, why it changed, and whether the original value is still recoverable.
\end{verification}
\end{labproject}
